\documentclass{article}
\usepackage[UTF8]{ctex}
\usepackage{amsmath,amssymb,amsfonts}
\usepackage{enumitem}
\usepackage{geometry}
\geometry{a4paper, margin=1in}
% === 错题本排版优化（不参与题目标识） ===
\setlength{\parindent}{0pt}
\setlength{\parskip}{0.35em plus 0.10em minus 0.08em}
\setlength{\emergencystretch}{3em}
\setlist{itemsep=0.25em, topsep=0.40em, parsep=0pt}
\allowdisplaybreaks[2]
% === 错题本排版优化结束 ===
\begin{document}
\section*{中山大学2024年数学分析真题}
\section*{152．中山大学2024年数学分析试卷}

\subsection{题目1}
1．（20 分）设 \(f(x)= \begin{cases}\frac{1}{x} e^{-\frac{1}{x^{2}}}, & x \neq 0 ; \\ a, & x=0 .\end{cases}\)

\begin{itemize}
  \item[（1）] 若 \(f(x)\) 在 \(x=0\) 连续，求 \(a\) 的值．
  \item[（2）] 在第（1）问的条件下，求 \(f^{\prime}(x)\) ．
  \item[（3）] \(f(x)\) 在 \(x=0\) 处是否可以进行泰勒展开？简要说明理由．
\end{itemize}

\subsection{题目2}
2．（20 分）证明如下结论：

\begin{itemize}
  \item[（1）] 已知 \(0<b<a<+\infty\) ，那么 \(\lim _{n \rightarrow \infty} \sqrt[n]{a^{n}+b^{n}}=a\) ．
  \item[（2）] 设 \(f(x)\) 在 \([a, b]\) 上连续，那么 \(\lim _{n \rightarrow \infty} \sqrt[n]{\int_{a}^{b}|f(x)|^{n} \mathrm{~d} x}=\max _{x \in[a, b]}|f(x)|\) ．
\end{itemize}

\subsection{题目3}
3．（10 分）已知 \(0<a<b<+\infty, f(x)\) 在 \([a, b]\) 上可导，求证：存在 \(\xi \in(a, b)\) ，使得

\[
\frac{a f(b)-b f(a)}{a-b}=f(\xi)-\xi f^{\prime}(\xi) .
\]

\subsection{题目4}
4．（10 分）已知函数列 \(\left\{f_{n}(x)\right\}\) 在区间 \(I\) 上一致收敛于 \(f(x)\) ，且对于每一个正整数 \(n, f_{n}(x)\) 均在 \(I\)上一致连续，证明：\(f(x)\) 在 \(I\) 上也一致连续．

\subsection{题目5}
5．（20 分）证明如下结论：

\begin{itemize}
  \item[（1）] 已知级数 \(\sum_{n=0}^{\infty} u_{n}\) 收敛，则 \(\lim _{n \rightarrow \infty} u_{n}=0\) ．
  \item[（2）] 举例若 \(f(x)\) 在 \([1,+\infty)\) 上连续且 \(\int_{1}^{+\infty} f(x) \mathrm{d} x\) 收敛，但 \(\lim _{x \rightarrow+\infty} f(x) \neq 0\) ．
  \item[（3）] 若 \(f(x)\) 在 \([1,+\infty)\) 上一致连续，加上（2）的条件，证明： \(\lim _{x \rightarrow+\infty} f(x)=0\) ．
\end{itemize}

\subsection{题目6}
6．（10 分）证明：\(f(x, y)=x y+\frac{1}{x}+\frac{1}{y}\) 在 \(0<x, y<+\infty\) 上有最小值，无最大值．

\subsection{题目7}
7．（10 分）求 \(y=\arctan x\) 和直线 \(x=1\) 及坐标 \(x\) 轴围成曲面三角形的面积．

\subsection{题目8}
8．（10 分）计算曲面积分 \(\iint_{\sum} x y z \mathrm{~d} x \mathrm{~d} y\) ，其中 \(\sum\) 为 \(x^{2}+y^{2}+z^{2}=1, x \geq 0, y \geq 0\) ，方向取外侧。

\subsection{题目9}
9．（20 分）设二元函数 \(u(x, y), v(x, y)\) 在区域 \(D=\left\{(x, y) \mid x^{2}+y^{2}<1\right\}\) 内连续可微，且 \(\frac{\partial u}{\partial y}=\frac{\partial v}{\partial x}\) ．

\begin{itemize}
  \item[（1）] 试求区域 \(D\) 上的二元函数 \(f(x, y)\)（用 \(u, v\) 来表示）满足 \(\frac{\partial f}{\partial x}=u, \frac{\partial f}{\partial y}=v\) 。
  \item[（2）] 若 \(D\) 换成 \(D^{\prime}=\left\{(x, y) \mid 0<x^{2}+y^{2}<1\right\}\) ，那么上述的函数 \(f(x, y)\) 是否仍总是存在？请说明理
\end{itemize}

\begin{itemize}
  \item[（1）] \(\int_{0}^{+\infty} e^{-x y} \sin y \mathrm{~d} y\) 在 \([0,+\infty)\) 上是否一致收敛？请说明理由。
  \item[（2）] \(\int_{0}^{+\infty} \mathrm{d} x \int_{0}^{+\infty}\left|e^{-x y} \sin y\right| \mathrm{d} y\) 是否有限？请说明理由。
  \item[（3）] 证明： \(\int_{0}^{+\infty} \mathrm{d} x \int_{0}^{+\infty} e^{-x y} \sin y \mathrm{~d} y=\int_{0}^{+\infty} \mathrm{d} x \int_{0}^{+\infty} e^{-x y} \sin x \mathrm{~d} y\) ．
\end{itemize}
\end{document}
